this shows that one obtains an isomorphism
\[
  A_0[t,t^{-1}]\longrightarrow A
\]
of graded \(A_0\)-algebras by sending \(t\) to \(f\), which completes
the proof of \SGARefLink{sga3:viii:corollary:4:7}{4.7}.

Under the hypotheses of
\SGARefLink{sga3:viii:proposition:4:6}{4.6},
\SGARefLink{sga3:viii:corollary:4:7}{4.7} already shows that the
\(\mathscr A_m\), for \(m\in M\), are invertible. To prove condition
\SGARefLink{sga3:viii:proposition:4:1:item:b:01}{4.1(b)}, we may
therefore suppose that \(\mathscr A_m\) and \(\mathscr A_{m'}\) have
bases \(f_m\) and \(f_{m'}\), with inverses
\[
  f_m^{-1}\in\Gamma(\mathscr A_{-m})
  \qquad\text{and}\qquad
  f_{m'}^{-1}\in\Gamma(\mathscr A_{-m'}).
\]
Multiplication by
\[
  f_m^{-1}f_{m'}^{-1}
  \in\Gamma(\mathscr A_{-m-m'})
\]
then defines a homomorphism
\[
  \mathscr A_{m+m'}\longrightarrow
  \mathscr A_0\simeq\mathcal O_S
\]
that sends the image of \(f_m\otimes f_{m'}\) to the unit section of
\(\mathcal O_S\). In the diagram
\SGARefTarget{sga3:viii:proposition:4:6:diagram:01}
\[
\begin{tikzcd}[column sep=large]
  \mathscr A_m\otimes\mathscr A_{m'}
    \arrow[r,"u"]
    \arrow[rr,bend left=32,"w"]
  & \mathscr A_{m+m'}
    \arrow[r,"v"]
  & \mathscr A_0\simeq\mathcal O_S ,
\end{tikzcd}
\]
the maps \(w\) and \(v\) are therefore epimorphisms of invertible
sheaves, hence isomorphisms; consequently \(u\) is an isomorphism.
\hfill\(\square\)
